% =====================================================================
%  A thermodynamically consistent framework of rate-independent
%  evolution: theory and variational approximation
%
%  Preprint version (journal-agnostic). To be reformatted with the
%  World Scientific ws-m3as class upon submission to
%  Mathematical Models and Methods in Applied Sciences (M3AS).
%
%  J. Ayensa-Jimenez, I. Romero
% =====================================================================
\documentclass[11pt,a4paper]{article}

% ---- encoding & fonts ----
\usepackage[utf8]{inputenc}
\usepackage[T1]{fontenc}
\usepackage{lmodern}

% ---- math ----
\usepackage{amsmath,amssymb,amsthm,mathtools}
\usepackage{bm}

% ---- layout ----
\usepackage[margin=2.6cm]{geometry}
\usepackage{microtype}
\usepackage{booktabs}
\usepackage{enumitem}
\usepackage{subcaption}
\usepackage{caption}

% ---- graphics: native pgfplots (no external binaries) ----
\usepackage{pgfplots}
\pgfplotsset{compat=1.18}
\usepgfplotslibrary{groupplots}

% ---- references ----
\usepackage[colorlinks=true,allcolors=blue!55!black]{hyperref}
\usepackage[capitalize,noabbrev]{cleveref}

% =====================================================================
%  Theorem environments
% =====================================================================
\theoremstyle{plain}
\newtheorem{theorem}{Theorem}[section]
\newtheorem{proposition}[theorem]{Proposition}
\newtheorem{lemma}[theorem]{Lemma}
\newtheorem{corollary}[theorem]{Corollary}

\theoremstyle{definition}
\newtheorem{definition}[theorem]{Definition}
\newtheorem{assumption}[theorem]{Assumption}
\newtheorem{remark}[theorem]{Remark}

% =====================================================================
%  Macros
% =====================================================================
\newcommand{\R}{\mathbb{R}}
\newcommand{\N}{\mathbb{N}}
\newcommand{\Q}{\mathcal{Q}}
\newcommand{\E}{\mathcal{E}}
\newcommand{\Diss}{\mathrm{Diss}}
\newcommand{\Dr}{\mathcal{D}}
\newcommand{\Var}{\mathrm{Var}}
\newcommand{\bs}[1]{\bm{#1}}
\newcommand{\dd}{\mathop{}\!\mathrm{d}}
\newcommand{\sgn}{\operatorname{sgn}}
\newcommand{\Id}{\mathrm{I}}
\DeclareMathOperator*{\argmin}{arg\,min}
\DeclareMathOperator{\dist}{dist}

\crefname{assumption}{Assumption}{Assumptions}
\Crefname{assumption}{Assumption}{Assumptions}

% =====================================================================
\begin{document}

\title{\bfseries A thermodynamically consistent framework of
rate-independent evolution:\\ theory and variational approximation}

\author{Jacobo Ayensa-Jim\'enez\thanks{Technical University of Madrid
(UPM), Madrid, Spain. \texttt{jacobo.ayensa@upm.es}.}
\and
Ignacio Romero\thanks{Technical University of Madrid (UPM) and IMDEA
Materials Institute, Madrid, Spain. \texttt{ignacio.romero@upm.es}.}}

\date{\today}

\maketitle

\begin{abstract}
We develop and analyse an abstract framework for the evolution of
dissipative systems within the theory of rate-independent processes.
The framework is specified by a triple consisting of a state space, a
stored-energy functional depending on the state and on a
time-dependent loading, and a $1$-homogeneous dissipation potential.
Assuming that the system evolves according to an energetic principle,
the governing equations---an energetic (global stability plus energy
balance) formulation and an equivalent subdifferential inclusion of
Biot type---are derived systematically, without additional modelling
hypotheses. We show that solutions satisfy, a priori, energetic
estimates intimately related to the first and second laws of
thermodynamics: the energy balance is an exact statement of energy
conservation, while the $1$-homogeneity of the dissipation potential
forces a non-negative, minimal (\emph{economical}) dissipation.
Under natural coercivity and continuity assumptions we establish the
existence of energetic solutions and, through a vanishing-viscosity
selection, of balanced-viscosity solutions that resolve the ambiguity
of the energetic formulation at jumps; uniqueness is obtained under
uniform convexity and, in the scalar case, under a mild
finite-multiplicity condition. We then construct a variational time
integrator based on incremental minimisation, prove discrete
stability and a discrete energy inequality, and establish convergence
to energetic solutions together with first- and second-order local
consistency estimates and a global energy-consistency bound.
The theory is illustrated on a scalar double-well benchmark driven by
a bounded, saturating loading, for which the return-map integrator
reproduces the exact quasi-static branches and exhibits the predicted
first-order energy consistency. The framework is motivated by, and
tailored to, a companion application to epigenetic evolution, developed
elsewhere.

\medskip
\noindent\textbf{Keywords:} rate-independent systems; energetic
solutions; balanced-viscosity solutions; dissipation potential;
hysteresis; variational time integration; return mapping;
thermodynamic consistency.

\medskip
\noindent\textbf{AMS Subject Classification:}
49J40, 74C05, 34G25, 74N30, 65M12, 80A17.
\end{abstract}

\tableofcontents
\bigskip

% =====================================================================
\section{Introduction}\label{sec:intro}
% =====================================================================

Many systems in the applied sciences are governed by evolution
equations that encapsulate the balance between energy storage and
dissipation. This kind of dissipative dynamics underlies, for
instance, the theory of plasticity~\cite{lubliner:1990,hanreddy2013},
viscoelasticity~\cite{tschoegl1989}, damage and
delamination~\cite{lemaitre2005uh}, diffusion in porous
media~\cite{bear1972}, frictional contact~\cite{wriggers2006},
phase transformation in polycrystals~\cite{balljames1987,mielke2002}
and ferromagnetism~\cite{desimone2002,brokate1996}. In all of these
settings the governing equations account for the energetic mechanisms
that take place during the process and ensure that fluxes balance the
accumulation of energy. Remarkably, thermodynamics provides a common
abstract structure that simplifies their analysis and reveals the
essential ingredients of a consistent theory~\cite{germain1983,
halphen1975generalized}.

A particularly rich subclass consists of \emph{rate-independent}
phenomena, whose behaviour is invariant under time
reparametrisation. We refer to the monographs of Mielke and
Roub{\'i}{\v c}ek~\cite{MielkeRoubicek}, of Visintin~\cite{visintin1994}
and of Krasnosel'ski{\u\i} and Pokrovski{\u\i}~\cite{Krasnoselskii}, and
to the foundational work of Mielke and Theil~\cite{MielkeTheil}, for
comprehensive accounts of the mathematical structure of such problems. As shown there,
rate-independent systems possess a variational structure that leads to
powerful analytical results and robust approximation schemes, valid
even in non-smooth settings. Using rate-independent, dissipative model
equations as a template, it becomes easier to propose new models,
irrespective of the field of application, knowing beforehand that the
resulting governing equations respect the most fundamental balance
principles.

Rate-independent models have proven extremely successful in mechanics
and, in particular, in elastoplasticity, fracture and damage
theories~\cite{francfort1998,dalmaso2006,carstensen2002uh}. However,
the framework is far more general, and its transfer to other
disciplines may open fruitful avenues for modelling and analysis.
The present work is motivated by one such transfer: a companion
paper develops a rate-independent model of \emph{epigenetic}
evolution, in which a scalar internal variable encodes the state of
chromatin between a healthy and a damaged phenotype, and an external
oxygen signal plays the role of the loading. Epigenetic marks such as
DNA methylation are heritable, long-lived, and yet actively
reversible~\cite{Bird2002}, and their dependence on the cellular
microenvironment---for example, the impairment of oxygen-dependent
demethylation enzymes under hypoxia~\cite{Thienpont2016}---produces
memory and hysteresis phenomena that are naturally captured by
rate-independent evolution. The biological development is deferred to
the companion paper; here we concentrate on the mathematical framework
that underpins it, which is of independent interest and applicable well
beyond the motivating example.

\paragraph{Scope and contributions.}
Following the abstract set-up of~\cite{MielkeRoubicek}, we specify a
rate-independent system by a triple $(\Q, E, \Psi)$ consisting of a
state space $\Q\subseteq\R^n$, a stored-energy functional $E$
depending on the state and on a prescribed loading path $S(t)$, and a
convex, $1$-homogeneous dissipation potential $\Psi$. The
contributions of this paper are the following.
\begin{enumerate}[leftmargin=2em,itemsep=2pt]
  \item We show that, once the triple is fixed and an energetic
  principle is postulated, the governing equations are determined
  systematically: we derive both the \emph{energetic formulation}
  (global stability and energy balance,
  \cref{def:energetic}) and the equivalent \emph{subdifferential} or
  Biot inclusion (\cref{sec:subdiff}), and we make precise the
  role of the $1$-homogeneity of $\Psi$ through a saturation identity
  (\cref{lem:saturation}).
  \item We give a self-contained thermodynamic reading of the
  framework (\cref{sec:thermo}): the energy balance is the first law,
  the non-negativity of the dissipation rate is the second law, and
  the saturation identity expresses a principle of \emph{minimal}
  (economical) dissipation. Over closed loading cycles this yields a
  purely geometric characterisation of the hysteretic dissipation.
  \item We collect and, where necessary, complete the well-posedness
  theory (\cref{sec:math}): properties of the induced dissipation
  distance, existence of energetic solutions, uniqueness under uniform
  convexity, and existence and (scalar) uniqueness of
  balanced-viscosity solutions obtained by vanishing viscosity, which
  resolve the non-uniqueness of the energetic formulation at jumps.
  \item We analyse a variational time integrator based on incremental
  minimisation (\cref{sec:numerical}): we prove discrete stability and
  a discrete energy inequality, convergence to energetic solutions,
  local consistency of orders one and two under weak and strong
  regularity, and a global energy-consistency estimate. In the scalar
  case the incremental problem reduces to a two-sided return map.
  \item We validate the theory on an abstract scalar double-well
  benchmark driven by a bounded, saturating loading
  (\cref{sec:experiments}), reproducing the analytically known
  quasi-static branches, the reversible and catastrophic hysteresis
  loops, and the predicted first-order energy consistency.
\end{enumerate}

\paragraph{Outline.}
\Cref{sec:framework} introduces the framework and the governing
equations. \Cref{sec:thermo} develops the thermodynamic
interpretation. \Cref{sec:math} presents the well-posedness results,
and \cref{sec:numerical} the variational numerical method and its
analysis. \Cref{sec:experiments} reports the numerical experiments,
and \cref{sec:conclusions} draws conclusions and outlines future work.
To keep the exposition readable, all proofs are collected in the
appendices.

% =====================================================================
\section{The rate-independent framework}\label{sec:framework}
% =====================================================================

\subsection{The defining triple}\label{sec:ingredients}

Following~\cite{MielkeRoubicek}, a rate-independent system is specified
by a triple $(\Q, E, \Psi)$.

\begin{enumerate}[label=(\roman*),leftmargin=2.2em,itemsep=3pt]
  \item \textbf{State space.} The set of \emph{admissible states} is a
  closed and convex subset of a Banach space; throughout we take
  \begin{equation}\label{eq:state_space}
    \Q \subseteq \R^n .
  \end{equation}

  \item \textbf{Loading.} The interaction of the system with its
  environment is encoded by a prescribed, time-dependent
  \emph{loading path}
  \begin{equation}
    S \in W^{1,1}(0,T;\R^{+}) .
  \end{equation}

  \item \textbf{Stored energy.} The \emph{stored-energy functional}
  $E\colon \Q\times\R^{+}\to\R$ is assumed to admit the additive
  decomposition
  \begin{equation}\label{eq:energy}
    E(\bs q, S) \;=\; W(\bs q)\;-\;\bs q\cdot\bs\ell(S),
  \end{equation}
  where $W\colon\Q\to\R$ is the \emph{configurational free energy} and
  $\bs\ell\colon\R^{+}\to\R^n$ the \emph{interaction (loading)
  potential}.

  \item \textbf{Dissipation potential.} The \emph{dissipation
  potential} is a convex, lower semicontinuous map
  $\Psi\colon\R^n\to\R$ satisfying
  \begin{itemize}[leftmargin=1.4em,itemsep=1pt]
    \item \emph{normalisation}: $\Psi(\bs 0)=0$;
    \item \emph{positiveness}: $\Psi(\bs v)\ge 0$ for all $\bs v\in\R^n$;
    \item \emph{$1$-homogeneity}:
    $\Psi(\lambda\bs v)=\lambda\,\Psi(\bs v)$ for all $\lambda>0$.
  \end{itemize}
  These are the structural assumptions required
  in~\cite{Mielke2005,MielkeRoubicek} for the dissipation potential of
  a rate-independent system; the $1$-homogeneity is the property that
  distinguishes rate-independent from viscous (rate-dependent)
  evolution. A map satisfying only the first two conditions is called
  a dissipation \emph{pseudo-potential}.
\end{enumerate}

The forthcoming well-posedness results require, in addition, the
quantitative hypotheses of \cref{ass:standing}, which are collected in
\cref{sec:math}. The benchmark of \cref{sec:experiments} satisfies all
of them.

\subsection{Derived canonical quantities}\label{sec:derived}

From the triple $(\Q,E,\Psi)$ the following objects are canonically
constructed.

\paragraph{(a) Thermodynamic driving force.}
When $E$ is differentiable in $\bs q$, the negative energy gradient
defines the \emph{generalised thermodynamic force}
\begin{equation}\label{eq:force}
  \bs f(\bs q, S) \;:=\; -\frac{\partial E}{\partial \bs q}(\bs q, S)
  \;=\; \bs\ell(S) - \bs g(\bs q),
  \qquad
  \bs g(\bs q) := \frac{\partial W}{\partial\bs q}(\bs q).
\end{equation}
The force $\bs f$ is the variable conjugate to $\bs q$: it quantifies
the thermodynamic drive toward evolution of the state.

\paragraph{(b) Elastic domain.}
The subdifferential of $\Psi$ at the origin,
\begin{equation}\label{eq:elastic_domain}
  \E := \partial\Psi(\bs 0),
\end{equation}
is the closed convex set of forces the system can sustain without
evolving. It is the \emph{elastic domain} in force space; by convex
duality~\cite{rockafellar1970,ekeland1999convex} $\Psi$ is its support
function, $\Psi(\bs v)=\sup_{\bs f\in\E}\bs f\cdot\bs v$.

\paragraph{(c) Dissipation distance.}
The \emph{dissipation distance} is the minimal cost of a transition
between two states,
\begin{equation}\label{eq:diss_dist}
  \Dr(\bs q_1,\bs q_2)
  := \inf\Bigl\{\textstyle\int_0^1 \Psi(\dot\gamma(s))\dd s :
  \gamma\in\mathrm{AC}([0,1];\Q),\ \gamma(0)=\bs q_1,\ \gamma(1)=\bs q_2\Bigr\}.
\end{equation}

\paragraph{(d) Total dissipation.}
For a trajectory $\bs q\colon[0,T]\to\Q$ of bounded variation, the
\emph{total dissipation} is
\begin{equation}\label{eq:total_diss}
  \Diss_\Psi(\bs q;[0,t])
  := \sup\Bigl\{\textstyle\sum_{j=1}^{N}\Dr\bigl(\bs q(t_{j-1}),\bs q(t_j)\bigr):
  0=t_0<\dots<t_N=t,\ N\in\N\Bigr\}.
\end{equation}

\paragraph{(e) Power and work of the loading.}
The rate of energy change due to the loading at frozen state is
\begin{equation}\label{eq:power}
  \mathcal P_S(t) := \partial_t E(\bs q, S(t))
  = \partial_S E(\bs q, S)\,\dot S(t)
  = -\bs q\cdot\bs\ell'(S(t))\,\dot S(t),
\end{equation}
and the associated \emph{work} performed on the system over $[0,t]$ is
$\mathcal W_S(t):=\int_0^t\partial_S E(\bs q(s),S(s))\,\dot S(s)\dd s$.

\subsection{Governing equations: energetic formulation}
\label{sec:energetic}

We postulate the evolution through the following variational notion,
due to Mielke and Theil~\cite{MielkeTheil}.

\begin{definition}[Energetic solution]\label{def:energetic}
A function $\bs q\colon[0,T]\to\Q$ is an \emph{energetic solution} of
$(\Q,E,\Psi)$ for the loading $S$ if it satisfies:

\smallskip
\noindent\textbf{(S) Global stability.} For every $t\in[0,T]$ and every
$\tilde{\bs q}\in\Q$,
\begin{equation}\label{eq:stability}
  E(\bs q(t),S(t)) \;\le\; E(\tilde{\bs q},S(t)) + \Dr(\bs q(t),\tilde{\bs q}).
\end{equation}

\noindent\textbf{(E) Energy balance.} For every $t\in[0,T]$,
\begin{equation}\label{eq:energy_balance}
  E(\bs q(t),S(t)) + \Diss_\Psi(\bs q;[0,t])
  = E(\bs q(0),S(0)) + \int_0^t \partial_S E(\bs q(s),S(s))\,\dot S(s)\dd s.
\end{equation}
\end{definition}

Condition~(S) admits a variational reading: $\bs q(t)$ is a global
minimiser of the perturbed energy $\bs r\mapsto E(\bs r,S(t))+\Dr(\bs
q(t),\bs r)$. Condition~(E) states that, along the trajectory, the sum
of stored energy and cumulative dissipation equals the initial energy
plus the work supplied by the loading; it is an exact energy
conservation statement, examined thermodynamically in
\cref{sec:thermo}.

\subsection{Governing equations: subdifferential formulation}
\label{sec:subdiff}

Since $\Psi$ is convex but generally non-smooth at $\dot{\bs q}=\bs 0$,
the pointwise force balance takes the form of a differential inclusion.

\begin{proposition}[Biot inclusion]\label{prop:biot}
Let $\bs q\in\mathrm{AC}([0,T];\Q)$ satisfy \textnormal{(S)}
and~\textnormal{(E)}, with $E(\cdot,S)$ differentiable. Then, for a.e.\
$t\in[0,T]$,
\begin{equation}\label{eq:biot}
  \bs 0 \;\in\; \partial\Psi(\dot{\bs q}(t)) + \frac{\partial E}{\partial\bs q}(\bs q(t),S(t)),
\end{equation}
equivalently $\bs f(\bs q,S)\in\partial\Psi(\dot{\bs q})$. Conversely,
if $E(\cdot,S)$ is convex, then~\eqref{eq:biot} together with the
local stability $\bs f(\bs q(t),S(t))\in\E$ implies
\textnormal{(S)}--\textnormal{(E)}.
\end{proposition}

Equation~\eqref{eq:biot} is the (generalised) Biot
equation~\cite{Biot1956}; for $\Psi=\rho\|\cdot\|$ it coincides with the
sweeping process of Moreau~\cite{moreau1977} and, in one dimension, with
the play operator of hysteresis~\cite{Krasnoselskii,visintin1994}. Its
analysis rests on the \emph{saturation identity}, the algebraic
signature of $1$-homogeneity: if $\bs f\in\partial\Psi(\bs v)$ then
$\bs f\cdot\bs v=\Psi(\bs v)$ (\cref{lem:saturation} in
\cref{app:framework}). Combined with \cref{prop:biot}, it yields, along
any solution,
\begin{equation}\label{eq:saturation_solution}
  \bs f(\bs q,S)\cdot\dot{\bs q}=\Psi(\dot{\bs q}),
\end{equation}
i.e.\ the dissipation is exactly balanced by the energy release at
frozen loading. This is the cornerstone of the thermodynamic
interpretation of \cref{sec:thermo}. The proof of \cref{prop:biot} is
given in \cref{app:framework}.

% =====================================================================
\section{Thermodynamic interpretation}\label{sec:thermo}
% =====================================================================

We now read the framework of \cref{sec:framework} in the language of
the thermodynamics of internal variables, identifying the energy
balance with the first law and the dissipation inequality with the
second law~\cite{germain1983,halphen1975generalized}.

\subsection{Free energy and the first law}\label{sec:first_law}

The state of the system is the pair $(\bs q, S)$, and $E$ plays the role
of a \emph{Helmholtz-type free energy}: the portion of the energy
stored in the configuration and recoverable through quasi-static
changes,
\begin{equation}\label{eq:free_energy}
  E(\bs q, S) = \underbrace{F(\bs q)}_{\text{configurational}}
  \;-\;\underbrace{\bs q\cdot\bs\ell(S)}_{\text{coupling to environment}}.
\end{equation}
The configurational term $F$ is intrinsic and independent of the
environment; the coupling term $-\bs q\cdot\bs\ell(S)$ acts as a
generalised chemical potential that biases the state. The conjugate
driving force is $\bs f=-\partial_{\bs q}E$, as in~\eqref{eq:force}, and
the power injected by the loading at frozen state is
$\mathcal P_S=\partial_S E\,\dot S$, as in~\eqref{eq:power}.

Rearranging the energy balance~\eqref{eq:energy_balance} over $[0,T]$
yields the \emph{first law}
\begin{equation}\label{eq:first_law}
  \underbrace{\Delta E}_{\text{change in free energy}}
  \;+\;
  \underbrace{\Diss_\Psi(\bs q;[0,T])}_{\mathcal D_{\mathrm{tot}}\ \ge 0}
  \;=\;
  \underbrace{\int_0^T \partial_S E(\bs q,S)\,\dot S\dd t}_{\mathcal W\ \text{(work of the environment)}},
\end{equation}
with $\Delta E=E(\bs q(T),S(T))-E(\bs q(0),S(0))$: the work supplied by
the environment is stored as free energy or irreversibly dissipated.

\subsection{Entropy production and the second law}\label{sec:second_law}

The instantaneous \emph{dissipation rate} is
\begin{equation}\label{eq:diss_rate}
  \mathcal D(t) := \bs f(\bs q(t),S(t))\cdot\dot{\bs q}(t).
\end{equation}

\begin{proposition}[Second law]\label{prop:second_law}
Along any solution of~\eqref{eq:biot} one has $\mathcal D(t)\ge 0$ for
a.e.\ $t$. Moreover the \emph{excess dissipation}
$\Xi(t):=\bs f(t)\cdot\dot{\bs q}(t)-\Psi(\dot{\bs q}(t))$ vanishes
identically, $\Xi\equiv 0$.
\end{proposition}

This is an immediate consequence of the saturation
identity~\eqref{eq:saturation_solution} (see \cref{app:framework}).
The non-negativity $\mathcal D\ge 0$ is the second law: dissipation is
non-negative regardless of the direction of evolution. The identity
$\Xi\equiv 0$ expresses a \emph{principle of minimal (economical)
dissipation}: a rate-independent system saturates the dissipation
inequality, dissipating exactly the minimum compatible with the energy
balance, with no viscous excess. In an isothermal setting at absolute
temperature $\Theta$, the irreversible entropy production rate is
$\Theta\,\dot s_{\mathrm{irr}}=\mathcal D(t)\ge0$, so that over the
whole process
\begin{equation}\label{eq:entropy}
  \Delta s_{\mathrm{irr}} = \frac{\mathcal D_{\mathrm{tot}}}{\Theta}
  = \frac{\Diss_\Psi(\bs q;[0,T])}{\Theta}\ \ge\ 0,
\end{equation}
the standard form of the second law.

\subsection{Geometric dissipation over closed cycles}\label{sec:hysteresis}

A hallmark of rate-independent thermodynamics is that the energy
dissipated over a closed loading cycle depends only on the geometry of
the trajectory, not on its time scale. Consider a cycle with
$S(0)=S(T)=S_0$ that returns the state to its initial value, so that
$\Delta E=0$. Integrating~\eqref{eq:first_law} and using
$\mathcal D\ge0$ (\cref{prop:second_law}) gives
\begin{equation}\label{eq:cycle}
  \mathcal W = \oint_0^T \partial_S E(\bs q,S)\,\dot S\dd t
  = \mathcal D_{\mathrm{tot}} \ \ge\ 0.
\end{equation}
Hence the work performed by the environment over a closed cycle cannot
be recovered and is entirely converted into irreversible entropy
production. In the $(\bs q,\bs f)$ plane this manifests as a
non-vanishing \emph{hysteresis loop} whose enclosed area equals exactly
$\mathcal D_{\mathrm{tot}}$; since $\bs f=\bs\ell(S)-\bs g(\bs q)$
(\cref{eq:force}), the same loop, up to the fixed reparametrisation
$\bs f\leftrightarrow\bs\ell(S)$, is the one plotted directly in the
$(\ell,q)$ loading--state plane in the numerical experiments of
\cref{sec:experiments}. The rate-independent structure therefore
guarantees that returning the environment to its baseline leaves a
permanent thermodynamic imprint---a property that, in the motivating
biological application, models the non-volatility of the epigenetic
state.

% =====================================================================
\section{Mathematical results}\label{sec:math}
% =====================================================================

Throughout this section we work under the following standing
hypotheses.

\begin{assumption}[Standing hypotheses]\label{ass:standing}
The triple $(\Q,E,\Psi)$ with $E(\bs q,S)=W(\bs q)-\bs q\cdot\bs\ell(S)$
satisfies:
\begin{enumerate}[label=\textnormal{(A\arabic*)},leftmargin=2.6em,itemsep=2pt]
  \item\label{A:W} \emph{Stored energy.} $\Q\subseteq\R^n$ is closed and
  $W\colon\Q\to\R$ is continuous and coercive of order $p>1$: there
  exist $c_1>0$, $c_2\ge0$ with $W(\bs q)\ge c_1\|\bs q\|^p-c_2$.
  \item\label{A:ell} \emph{Interaction potential.}
  $\bs\ell\colon[0,\infty)\to\R^n$ is Lipschitz and bounded, with
  $C_\ell:=\sup_{S\ge0}\|\bs\ell(S)\|<\infty$ and
  $C_{\ell'}:=\sup_{S\ge0}\|\bs\ell'(S)\|<\infty$.
  \item\label{A:S} \emph{Loading.} $S\in W^{1,1}(0,T;\R^+)$.
  \item\label{A:Psi} \emph{Dissipation potential.}
  $\Psi\colon\R^n\to[0,\infty)$ is continuous, convex, normalised
  $(\Psi(\bs 0)=0)$, positively $1$-homogeneous and symmetric
  $(\Psi(-\bs v)=\Psi(\bs v))$.
  \item\label{A:coerc} \emph{Definiteness of dissipation.} There
  exists $c_\Psi>0$ with $\Psi(\bs v)\ge c_\Psi\|\bs v\|$ for all
  $\bs v$; equivalently, $\bs 0\in\operatorname{int}\E$.
\end{enumerate}
\end{assumption}

Hypothesis~\ref{A:coerc} is used only for the total-variation bound in
the balanced-viscosity theory (\cref{sec:bv}); it holds with
$c_\Psi=\rho$ for $\Psi(\bs v)=\rho\|\bs v\|$ and, more generally,
whenever the elastic domain has the origin in its interior.

\subsection{Properties of the dissipation distance}\label{sec:diss_props}

\begin{proposition}[Quasi-pseudodistance]\label{prop:quasi_distance}
Under~\ref{A:Psi}, the dissipation distance~\eqref{eq:diss_dist}
satisfies, for all $\bs q_1,\bs q_2,\bs q_3\in\Q$,
\[
  \Dr(\bs q_1,\bs q_2)\ge0,\quad \Dr(\bs q,\bs q)=0,\qquad
  \Dr(\bs q_1,\bs q_3)\le\Dr(\bs q_1,\bs q_2)+\Dr(\bs q_2,\bs q_3).
\]
If $\Psi$ is symmetric then $\Dr$ is a pseudodistance; if in addition
$\Psi(\bs v)=0\Rightarrow\bs v=\bs 0$ (e.g.\ under~\ref{A:coerc}), then
$\Dr$ is a distance.
\end{proposition}

\begin{proposition}[Linearity and continuity]\label{prop:Dr_linear}
If $\Q$ is convex and $\Psi$ is continuous, convex and positively
$1$-homogeneous, then
\begin{equation}\label{eq:Dr_linear}
  \Dr(\bs q_1,\bs q_2)=\Psi(\bs q_2-\bs q_1)\qquad\forall\,\bs q_1,\bs q_2\in\Q,
\end{equation}
and $\Dr$ is continuous on $\Q\times\Q$. In particular, for
$\Psi(\dot q)=\rho|\dot q|$ one has $\Dr(q_1,q_2)=\rho|q_2-q_1|$.
\end{proposition}

The identity~\eqref{eq:Dr_linear} is central to the numerical method:
it makes the incremental dissipation independent of the step size.

\subsection{Energy estimates}\label{sec:energy_estimates}

\begin{proposition}[Coercivity, absolute continuity and power
control]\label{prop:power}
Under~\ref{A:W}--\ref{A:S}, the following hold.
\begin{enumerate}[label=\textnormal{(E\arabic*)},leftmargin=2.6em,itemsep=2pt]
  \item\label{E:compact} For every $c\in\R$ and $S\ge0$ the sublevel
  set $\{\bs q\in\Q:E(\bs q,S)\le c\}$ is compact.
  \item\label{E:AC} For each $\bs q\in\Q$, $t\mapsto E(\bs q,S(t))$ is
  absolutely continuous with
  $\frac{\dd}{\dd t}E(\bs q,S(t))=-\bs q\cdot\bs\ell'(S(t))\dot S(t)$
  a.e.
  \item\label{E:gronwall} There exist $c_3,c_4>0$, depending only on
  the data, with $|\partial_S E(\bs q,S)|\le c_3\bigl(E(\bs q,S)+c_4\bigr)$
  for all $\bs q\in\Q$, $S\ge0$.
\end{enumerate}
\end{proposition}

\subsection{Existence of energetic solutions}\label{sec:existence}

We first recall the set of states that pass the stability test.

\begin{definition}[Stability set]\label{def:stability_set}
For $t\in[0,T]$, the \emph{stability set} is
$\mathcal S(t):=\{\bs q\in\Q: E(\bs q,S(t))\le E(\tilde{\bs q},S(t))+\Dr(\bs q,\tilde{\bs q})\ \ \forall\,\tilde{\bs q}\in\Q\}$.
\end{definition}

The passage to the limit in the existence proof rests on the following
compatibility properties, which we state explicitly (the corresponding
statement is only sketched in the general theory).

\begin{lemma}[Compatibility]\label{lem:compatibility}
Under~\ref{A:W}--\ref{A:Psi}, let $(t_m,\bs q_m)$ be a stable sequence,
i.e.\ $\bs q_m\in\mathcal S(t_m)$ with
$\sup_m E(\bs q_m,S(t_m))<\infty$, and suppose $t_m\to t$ and
$\bs q_m\to\bs q$ in $\Q$. Then:
\textnormal{(C1)} the power converges, $\partial_S E(\bs q_m,S(t_m))\to
\partial_S E(\bs q,S(t))$ for a.e.\ $t$; and \textnormal{(C2)} the
stability set is closed, $\bs q\in\mathcal S(t)$.
\end{lemma}

\begin{theorem}[Existence of energetic solutions]\label{thm:existence}
Under~\ref{A:W}--\ref{A:Psi}, for every initial datum
$\bs q_0\in\mathcal S(0)$ there exists at least one energetic solution
$\bs q\colon[0,T]\to\Q$ of bounded variation, in the sense of
\cref{def:energetic}.
\end{theorem}

The proof verifies the abstract hypotheses of
\cite[Thm.~2.1.6]{MielkeRoubicek} using
\cref{prop:quasi_distance,prop:power,lem:compatibility}; see
\cref{app:math}. For the benchmark of \cref{sec:experiments}, $\Q=\R$
is not compact but $W$ is a coercive quartic, so~\ref{A:W} holds and
\cref{thm:existence} applies.

\subsection{Uniqueness under uniform convexity}\label{sec:uniqueness}

Energetic solutions need not be unique: at a jump, several post-jump
states may satisfy~(S)--(E). Uniqueness is recovered under convexity.

\begin{theorem}[Uniqueness]\label{thm:uniqueness}
Assume, in addition to~\ref{A:W}--\ref{A:Psi}, that
$W\in C^3(\Q)$ is uniformly convex ($D^2W\succeq\alpha\Id$, $\alpha>0$)
and that $\bs\ell\in C^3$, $S\in C^3([0,T])$. Then for every
$\bs q_0\in\mathcal S(0)$ the energetic solution is unique.
\end{theorem}

Uniform convexity forces a single well and thus excludes the
multi-well energies of interest; the scalar theory of \cref{sec:bv}
provides an alternative route.

\subsection{Balanced-viscosity solutions}\label{sec:bv}

\paragraph{Motivation.}
At a jump point $t^\ast$, conditions~(S)--(E) require the released
energy to cover the dissipation cost
$\Dr(\bs q^-(t^\ast),\bs q^+(t^\ast))$ but do not select the post-jump
state. Balanced-viscosity (BV) solutions~\cite{mielke2012bv} resolve
this by selecting the jump path as the vanishing-viscosity limit
($\varepsilon\to0^+$) of the regularised inclusion
\begin{equation}\label{eq:regularised}
  \bs 0\in\partial\Psi(\dot{\bs q}_\varepsilon)
  +\varepsilon\,\dot{\bs q}_\varepsilon
  +\frac{\partial E}{\partial\bs q}(\bs q_\varepsilon,S),
\end{equation}
whose memory is retained in a correction concentrated at jumps.

\begin{definition}[Balanced-viscosity solution]\label{def:BV}
A function $\bs q\colon[0,T]\to\Q$ of bounded variation is a
\emph{BV solution} of $(\Q,E,\Psi)$ if:

\smallskip
\noindent\textbf{(BV-S) Local stability.} For a.e.\ $t$, there is
$\delta>0$ such that $E(\bs q(t),S(t))\le E(\bs r,S(t))+\Dr(\bs q(t),\bs r)$
for all $\bs r\in B_\delta(\bs q(t))\cap\Q$.

\smallskip
\noindent\textbf{(BV-E) Energy balance with viscous correction.}
For every $t$,
\begin{equation}\label{eq:BV_energy}
  E(\bs q(t),S(t))+\Diss_\Psi(\bs q;[0,t])+\mathcal V(\bs q;[0,t])
  = E(\bs q_0,S(0))+\int_0^t\partial_S E(\bs q,S)\dot S\dd s,
\end{equation}
where the \emph{viscous correction} is
\begin{equation}\label{eq:viscous_correction}
  \mathcal V(\bs q;[0,t]):=\sum_{s\le t}
  \bigl[E(\bs q^-(s),S(s))-E(\bs q^+(s),S(s))-\Dr(\bs q^-(s),\bs q^+(s))\bigr]^+,
\end{equation}
the sum ranging over the (at most countably many) jump points.
\end{definition}

\begin{remark}[Jump condition]\label{rem:jump}
$\mathcal V$ is fully determined by the jump structure: at each jump
$t^\ast$ the summand encodes the \emph{jump condition}
$E(\bs q^-(t^\ast),S(t^\ast))-E(\bs q^+(t^\ast),S(t^\ast))
=\Dr(\bs q^-(t^\ast),\bs q^+(t^\ast))+\text{(viscous excess)}$,
selecting the path of steepest descent through the energy barrier.
\end{remark}

\begin{theorem}[Existence of BV solutions]\label{thm:BV_existence}
Under~\ref{A:W}--\ref{A:coerc}, for every $\bs q_0\in\mathcal S(0)$
there exists at least one BV solution in the sense of \cref{def:BV}.
\end{theorem}

\begin{theorem}[Scalar uniqueness]\label{thm:BV_uniqueness}
Let $n=1$, $\Q\subseteq\R$, and assume, besides
\ref{A:W}--\ref{A:coerc}, that $q\mapsto E(q,S)$ has at most finitely
many local minima for each $S\ge0$. Then the BV solution is unique.
\end{theorem}

\cref{thm:BV_uniqueness} covers the scalar double-well energy of the
benchmark and of the biological application, for which
\cref{thm:uniqueness} fails. Since every energetic solution is a BV
solution, uniqueness of BV solutions implies that the energetic
solution, when it is BV, is the unique one. The proofs of the results
of this section are collected in \cref{app:math}.

% =====================================================================
\section{Variational numerical method}\label{sec:numerical}
% =====================================================================

\subsection{Incremental minimisation}\label{sec:incremental}

We approximate energetic solutions by incremental
minimisation~\cite{ortiz1999tq,mielke2008zb,conti2008,romero2021dd}. Partition
$[0,T]$ as $0=t_0<t_1<\dots<t_N=T$, write $h_k=t_{k+1}-t_k$,
$S_k=S(t_k)$, and, given $\bs q_k$, define
\begin{equation}\label{eq:Ik}
  \bs q_{k+1}\in\argmin_{\bs q\in\Q} I_k(\bs q),
  \qquad
  I_k(\bs q)=E(\bs q,S_{k+1})-E(\bs q_k,S_k)+h_k\,\Psi\!\Bigl(\tfrac{\bs q-\bs q_k}{h_k}\Bigr).
\end{equation}
By the $1$-homogeneity of $\Psi$ and \cref{prop:Dr_linear}, minimising
$I_k$ is equivalent to minimising
\begin{equation}\label{eq:Ikstar}
  I_k^\ast(\bs q)=E(\bs q,S_{k+1})+\Dr(\bs q_k,\bs q),
\end{equation}
whose optimality condition is the backward-Euler discretisation
of~\eqref{eq:biot},
\begin{equation}\label{eq:discrete_inclusion}
  \bs 0\in\frac{\partial E}{\partial\bs q}(\bs q_{k+1},S_{k+1})
  +\partial\Psi(\bs q_{k+1}-\bs q_k).
\end{equation}

\begin{proposition}[Return map]\label{prop:return_map}
Under~\ref{A:Psi}, the incremental problem~\eqref{eq:Ikstar} is
equivalent to~\eqref{eq:discrete_inclusion}, which splits into two
mutually exclusive cases:
\begin{enumerate}[label=\textnormal{(\roman*)},leftmargin=2.4em,itemsep=2pt]
  \item \emph{Elastic lock:} $\bs q_{k+1}=\bs q_k$ iff
  $-\partial_{\bs q}E(\bs q_k,S_{k+1})\in\E$;
  \item \emph{Dissipative flow:} $\bs q_{k+1}\ne\bs q_k$ iff
  $-\partial_{\bs q}E(\bs q_{k+1},S_{k+1})\in\partial\Psi(\bs q_{k+1}-\bs q_k)\setminus\E$.
\end{enumerate}
In the scalar case $\Psi(\dot q)=\rho|\dot q|$, with trial force
$f^{\mathrm{tr}}=\ell(S_{k+1})-g(q_k)$, the update reads
\begin{equation}\label{eq:scalar_return_map}
  q_{k+1}=
  \begin{cases}
    q_k, & |f^{\mathrm{tr}}|\le\rho \quad(\text{lock}),\\[2pt]
    \text{solution of } g(q_{k+1})=\ell(S_{k+1})-\sgn(f^{\mathrm{tr}})\,\rho,
      & |f^{\mathrm{tr}}|>\rho\quad(\text{flow}).
  \end{cases}
\end{equation}
\end{proposition}

Equation~\eqref{eq:scalar_return_map} is a two-sided return
map~\cite{Simo1986,hanreddy2013}: a trial (elastic) predictor followed,
if the yield surface $|f|=\rho$ is violated, by a plastic corrector
onto it.

\subsection{Discrete stability and energy inequality}\label{sec:discrete_stab}

\begin{proposition}[Discrete stability and energy
bound]\label{prop:discrete_energy}
Let $\{\bs q_k\}$ be generated by~\eqref{eq:Ikstar}. Then for each $k$,
\begin{align}
  &E(\bs q_{k+1},S_{k+1})\le E(\tilde{\bs q},S_{k+1})+\Dr(\bs q_k,\tilde{\bs q})
  \quad\forall\,\tilde{\bs q}\in\Q,
  \label{eq:disc_stab}\\
  &E(\bs q_{k+1},S_{k+1})+\Dr(\bs q_k,\bs q_{k+1})
  \le E(\bs q_k,S_{k+1})
  = E(\bs q_k,S_k)+\bigl[E(\bs q_k,S_{k+1})-E(\bs q_k,S_k)\bigr].
  \label{eq:disc_energy}
\end{align}
\end{proposition}

The bound~\eqref{eq:disc_energy} is generally a strict inequality: the
scheme may over-dissipate at the discrete level because the external
work is evaluated at the frozen state $\bs q_k$. As $h\to0$ the gap
telescopes to the exact continuous balance~\eqref{eq:energy_balance},
which is the content of the next result.

\subsection{Convergence}\label{sec:convergence}

\begin{theorem}[Convergence of the incremental scheme]\label{thm:convergence}
Under~\ref{A:W}--\ref{A:Psi}, let $\bs q_0\in\mathcal S(0)$ and let
$\bs q^h$ denote the piecewise-constant interpolant of the sequence
$\{\bs q_k\}$ generated by~\eqref{eq:Ikstar} on a partition of maximal
step $h$. Then, as $h\to0$, there is a subsequence $h_j\to0$ such that
$\bs q^{h_j}(t)\to\bs q(t)$ for every $t\in[0,T]$, where $\bs q$ is an
energetic solution of $(\Q,E,\Psi)$. If, moreover, the solution is
unique (\cref{thm:uniqueness} or~\cref{thm:BV_uniqueness}), the whole
family $\bs q^h$ converges.
\end{theorem}

\subsection{Consistency and error bounds}\label{sec:consistency}

Substituting the exact solution into the discrete energy
bound~\eqref{eq:disc_energy} produces, since the exact solution obeys
the continuous balance~\eqref{eq:energy_balance}, the \emph{local
truncation error}
\begin{equation}\label{eq:truncation}
  \tau_k := \int_{t_k}^{t_{k+1}}\!\partial_S E(\bs q(s),S(s))\,\dot S(s)\dd s
  -\bigl[E(\bs q(t_k),S_{k+1})-E(\bs q(t_k),S_k)\bigr],
\end{equation}
the discrepancy between the continuous work and its frozen-state
discrete approximation.

\begin{proposition}[Local consistency]\label{prop:consistency}
Assume~\ref{A:ell}, $S\in W^{1,\infty}(0,T)$, and let $\bs q$ be the
exact solution.
\begin{enumerate}[label=\textnormal{(\roman*)},leftmargin=2.4em,itemsep=2pt]
  \item \emph{Weak regularity.} If $\bs q\in\mathrm{AC}(0,T;\Q)$, then
  \begin{equation}\label{eq:cons_weak}
    |\tau_k|\le C_{\ell'}\|\dot S\|_{L^\infty}
    \Bigl(\int_{t_k}^{t_{k+1}}\|\dot{\bs q}\|\dd\tau\Bigr)h_k
    = O(h_k).
  \end{equation}
  \item \emph{Strong regularity.} If $\bs q\in W^{1,\infty}(0,T;\Q)$,
  then
  \begin{equation}\label{eq:cons_strong}
    |\tau_k|\le\tfrac12\,C_{\ell'}\|\dot{\bs q}\|_{L^\infty}\|\dot S\|_{L^\infty}\,h_k^2
    = O(h_k^2).
  \end{equation}
\end{enumerate}
\end{proposition}

\begin{theorem}[Global energy consistency]\label{thm:global_consistency}
Let $h=\max_k h_k$ and
$\mathcal E_\tau:=\sum_{k=0}^{N-1}|\tau_k|$. Under the hypotheses of
\cref{prop:consistency},
\begin{equation}\label{eq:global_cons}
  \mathcal E_\tau\le C\,h,\qquad
  C=\begin{cases}
    C_{\ell'}\|\dot S\|_{L^\infty}\int_0^T\|\dot{\bs q}\|\dd\tau, & \bs q\in\mathrm{AC},\\[4pt]
    \tfrac12\,C_{\ell'}\|\dot{\bs q}\|_{L^\infty}\|\dot S\|_{L^\infty}\,T, & \bs q\in W^{1,\infty},
  \end{cases}
\end{equation}
so the scheme is globally first-order consistent in the energy balance,
irrespective of the regularity of $\bs q$.
\end{theorem}

The regularity of $\bs q$ is inherited from the smoothness of the data
$W,\bs\ell,S$ away from jumps; at a jump $\bs q$ is only of bounded
variation and the local order degrades, but the global energy
consistency $\mathcal E_\tau=O(h)$ persists. All proofs of this section
are given in \cref{app:numerical}.

% =====================================================================
\section{Numerical experiments}\label{sec:experiments}
% =====================================================================

We illustrate the theory on a scalar double-well benchmark chosen so
that all hypotheses of \cref{ass:standing} hold and the exact
quasi-static branches are known in closed form.

\subsection{Benchmark}\label{sec:benchmark}

Let $\Q=\R$ and
\begin{equation}\label{eq:benchmark}
  W(q)=k\,q^2(1-q)^2,\qquad
  \ell(S)=\ell_\infty\bigl(1-e^{-\lambda S}\bigr),\qquad
  \Psi(\dot q)=\rho\,|\dot q|,
\end{equation}
with $k,\ell_\infty,\lambda,\rho>0$, so that
$E(q,S)=k\,q^2(1-q)^2-q\,\ell(S)$ and
$g(q)=W'(q)=2k\,q(1-q)(1-2q)$. The loading $\ell$ is strictly
increasing and saturating: a \emph{high} signal $S$ drives $q$ toward
the second well $q=1$. It is bounded ($C_\ell=\ell_\infty$) and
Lipschitz ($C_{\ell'}=\ell_\infty\lambda$), so~\ref{A:ell} holds;
$W$ is a coercive quartic, so~\ref{A:W} holds with $p=4$; and
$\Psi=\rho|\cdot|$ satisfies~\ref{A:Psi} and~\ref{A:coerc} with
$c_\Psi=\rho$. The energy $E(\cdot,S)$ has at most two local minima,
so \cref{thm:BV_uniqueness} guarantees a unique BV solution.

The lower branch of $g$ attains its maximum
\begin{equation}\label{eq:gmax}
  g_{\max}=\frac{k}{3\sqrt3}
  \qquad\text{at}\qquad q_-=\frac{3-\sqrt3}{6},
\end{equation}
the fold beyond which the near branch of $g(q)=\tau$ disappears and the
system jumps. On the yield surface $g(q)=\tau$ (with
$\tau=\ell-\sgn(f)\rho$), the admissible root on $[0,q_-]$ is given by
the trigonometric (Cardano) formula
\begin{equation}\label{eq:cubic}
  q=\tfrac12+\tfrac{1}{\sqrt3}\cos\!\Bigl(\tfrac13\arccos\!\bigl(\tfrac{\tau}{g_{\max}}\bigr)-\tfrac{4\pi}{3}\Bigr),
  \qquad |\tau|<g_{\max}.
\end{equation}
We drive the system with the smooth cyclic loading
$S(t)=\tfrac12 S_{\max}\bigl(1-\cos(2\pi t/T)\bigr)$, $t\in[0,T]$, which
increases from $0$ to $S_{\max}$ and back. We fix $k=1$, $\rho=0.03$,
$\lambda=1$, $S_{\max}=5$, $T=1$, and consider two regimes:
a \emph{reversible} one ($\ell_\infty=0.18$, so
$\ell_\infty-\rho=0.15<g_{\max}\approx0.1925$) and a
\emph{catastrophic} one ($\ell_\infty=0.30$, so
$\ell_\infty-\rho=0.27>g_{\max}$).

\subsection{Qualitative behaviour}\label{sec:qualitative}

\Cref{fig:landscape} shows how the loading tilts the double-well
landscape, lowering the barrier toward $q=1$ as $S$ grows.
\Cref{fig:evolution}(a) shows the reversible cycle: the state remains
locked (elastic) until the driving force reaches the yield threshold,
then flows, and on unloading locks again before partially recovering,
leaving a residual $q_{\mathrm{res}}>0$. \Cref{fig:evolution}(b) shows
the resulting hysteresis loops; in the catastrophic regime the near
branch disappears at the fold and the state jumps to the far well
($q\approx1$), the signature of an abrupt transition.

The closed-form predictions~\eqref{eq:cubic} agree with the computed
extrema to four significant digits, as reported in
\cref{tab:validation}.

\begin{figure}[t]
  \centering
  \begin{tikzpicture}
    \begin{axis}[width=0.62\textwidth,height=0.34\textwidth,
      xlabel=$q$,ylabel=$E(q,S)$,
      legend pos=north west,legend cell align=left,
      grid=both,grid style={gray!20},tick label style={font=\small}]
      \addplot[thick,blue] table[x=q,y=E]{figures/fig_landscape_lo.dat};
      \addplot[thick,teal!70!black] table[x=q,y=E]{figures/fig_landscape_mid.dat};
      \addplot[thick,red!70!black] table[x=q,y=E]{figures/fig_landscape_hi.dat};
      \legend{$S=0$,$S=1$,$S=5$}
    \end{axis}
  \end{tikzpicture}
  \caption{Tilting of the double-well energy $E(\cdot,S)$ as the loading
  $S$ increases (reversible regime, $\ell_\infty=0.18$). The right well
  is progressively favoured.}
  \label{fig:landscape}
\end{figure}

\begin{figure}[t]
  \centering
  \begin{subfigure}[t]{0.49\textwidth}
    \centering
    \begin{tikzpicture}
      \begin{axis}[width=\textwidth,height=0.72\textwidth,
        xmin=0,xmax=1,axis y line*=left,xlabel=$t$,ylabel=$q(t)$,
        ymin=-0.01,tick label style={font=\small}]
        \addplot[thick,blue] table[x=t,y=q]{figures/fig_time_smooth.dat};
      \end{axis}
      \begin{axis}[width=\textwidth,height=0.72\textwidth,
        xmin=0,xmax=1,axis y line*=right,axis x line=none,ylabel=$S(t)$,
        tick label style={font=\small},legend pos=north east,
        legend cell align=left]
        \addlegendimage{thick,blue}\addlegendentry{$q(t)$}
        \addplot[thick,red!70!black,dashed] table[x=t,y=S]{figures/fig_time_smooth.dat};
        \addlegendentry{$S(t)$}
      \end{axis}
    \end{tikzpicture}
    \caption{Time evolution (reversible)}
  \end{subfigure}
  \hfill
  \begin{subfigure}[t]{0.49\textwidth}
    \centering
    \begin{tikzpicture}
      \begin{axis}[width=\textwidth,height=0.72\textwidth,
        xlabel=$\ell(S)$,ylabel=$q$,legend pos=north west,
        legend cell align=left,grid=both,grid style={gray!20},
        tick label style={font=\small}]
        \addplot[thick,blue] table[x=ell,y=q]{figures/fig_loop_smooth.dat};
        \addplot[thick,red!70!black] table[x=ell,y=q]{figures/fig_loop_cat.dat};
        \legend{reversible,catastrophic}
      \end{axis}
    \end{tikzpicture}
    \caption{Hysteresis loops}
  \end{subfigure}
  \caption{(a) Elastic lock, dissipative flow and partial recovery over
  a loading cycle. (b) Hysteresis loops in the $(\ell,q)$ plane: the
  reversible loop stays on the near branch, while the catastrophic loop
  jumps at the fold $g_{\max}$ to the far well.}
  \label{fig:evolution}
\end{figure}

\begin{table}[t]
  \centering
  \caption{Analytical~\eqref{eq:cubic} vs.\ computed extrema
  (reversible regime).}
  \label{tab:validation}
  \begin{tabular}{lccc}
    \toprule
    & analytical & numerical & rel.\ error \\
    \midrule
    $q_{\max}$ (peak damage)   & $0.10540$ & $0.10535$ & $5\cdot10^{-4}$ \\
    $q_{\mathrm{res}}$ (residual scar) & $0.01578$ & $0.01573$ & $3\cdot10^{-3}$ \\
    \bottomrule
  \end{tabular}
\end{table}

\subsection{Convergence and consistency}\label{sec:conv_exp}

We verify the consistency theory of \cref{sec:consistency} on the
reversible cycle. Two error measures are computed against a fine
reference ($N=2\cdot10^5$): the $L^\infty$ error of the state
$\|q^h-q_{\mathrm{ref}}\|_\infty$, and the global energy-balance defect
$\sum_k e_k=(E_0+\mathcal W_h)-(E_N+\Diss_h)\ge0$, the accumulated gap
in the discrete energy inequality~\eqref{eq:disc_energy}.

Two facts emerge, both consistent with the theory. First, the state
error is at the level of machine precision ($\sim10^{-16}$): during
smooth flow the return map~\eqref{eq:scalar_return_map} solves the
algebraic yield condition $g(q_{k+1})=\ell(S_{k+1})-\sgn(f)\rho$
\emph{exactly} at each node, independently of $h$, so the incremental
scheme reproduces the exact quasi-static branch at the grid points.
The entire discretisation error is thus carried by the energy balance,
not by the state. Second, the energy defect decays as $O(h)$
(\cref{fig:convergence}), exactly the first-order global energy
consistency of \cref{thm:global_consistency}: halving $h$ halves the
defect, and the data track the reference slope-one line.

\begin{figure}[t]
  \centering
  \begin{tikzpicture}
    \begin{loglogaxis}[width=0.6\textwidth,height=0.4\textwidth,
      xlabel=$h$,ylabel=$\sum_k e_k$ (energy defect),
      legend pos=north west,legend cell align=left,
      grid=both,grid style={gray!20},tick label style={font=\small}]
      \addplot[only marks,mark=*,blue] table[x=h,y=energy_defect]{figures/fig_convergence.dat};
      \addplot[dashed,black,domain=1.5e-4:2e-2]{0.0748*x};
      \legend{energy defect,slope $1$}
    \end{loglogaxis}
  \end{tikzpicture}
  \caption{Global energy-balance defect versus step size $h$
  (log--log). The first-order decay confirms
  \cref{thm:global_consistency}. The state error, by contrast, is at
  machine precision because the return map is exact on smooth branches.}
  \label{fig:convergence}
\end{figure}

% =====================================================================
\section{Conclusions and future work}\label{sec:conclusions}
% =====================================================================

We have proposed Rate-Independent Epigenetics (RIE), a framework that
models epigenetic change as rate-independent dissipative evolution
driven by a single triple $(\Q,E,\Psi)$: an energy landscape of
chromatin configurations, a micro-environmental loading, and a
$1$-homogeneous dissipation potential encoding the resistance to
remodelling. Once the triple is fixed and an energetic principle is
postulated, the governing equations---the energetic formulation and the
equivalent Biot inclusion---follow systematically, and their solutions
satisfy, a priori, estimates that coincide with the first and second
laws of thermodynamics. The $1$-homogeneity of the dissipation
potential is the structural source of rate independence: it yields the
saturation identity, the non-negativity and minimality of the
dissipation, and the purely geometric characterisation of the
hysteretic dissipation over closed cycles. Threshold activation,
memory, and residual state upon unloading---the defining phenomenology
of epigenetic marks---are thus not modelled by ad-hoc equations but
emerge from a thermodynamically consistent variational structure.

On the analytical side we established the well-posedness of RIE models,
making explicit the compatibility argument behind existence and adding
the definiteness hypothesis needed for the balanced-viscosity
construction. Energetic solutions exist under coercivity and
continuity; they are unique under convexity and, for the scalar
bistable landscapes of epigenetic interest, the balanced-viscosity
selection restores uniqueness without convexity. On the numerical side
we analysed an incremental variational integrator: it is
unconditionally stable in the energetic sense, converges to energetic
solutions, and is first-order globally consistent in the energy
balance, reducing in the scalar case to a two-sided return map. The
linear play operator---the minimal RIE model of epigenetic
memory---admits a closed-form evolution exhibiting both reversible and
irreversible responses, and the integrator reproduced it exactly at the
nodes, confining the discretisation error to the energy balance.

Several directions remain open. The balanced-viscosity theory is
developed here for scalar states; its extension to vectorial
epigenetic variables, where the transition path at a switch is
genuinely path-dependent, is an important next step, as is a rigorous
convergence rate for the state in the presence of switches. The
local truncation estimate invites adaptive time-stepping, and
non-symmetric dissipation potentials would model the
direction-dependent resistance of methylation versus demethylation.

Most importantly, the present framework is the abstract foundation for
a quantitative epigenetic model. In a companion paper we specialise RIE
to a scalar \emph{bistable} landscape---a double-well free energy with a
saturating micro-environmental coupling---in which the healthy and
damaged phenotypes are the two wells. There the convexity of the
present benchmark is lost, catastrophic phenotype switches (of
epithelial--mesenchymal type) appear as energetic jumps selected by the
balanced-viscosity criterion, and the analytical thresholds, residual
``epigenetic scars'' and bifurcation structure can be calibrated
against experimental data on hypoxia-driven epigenetic remodelling.


\appendix
% =====================================================================
%  Appendices — proofs
% =====================================================================

\section{Proofs for the rate-independent framework}
\label{app:framework}

\begin{proof}[Proof of \cref{lem:saturation}]
Let $\bs f\in\partial\Psi(\bs v)$. By the subgradient inequality,
$\Psi(\bs w)\ge\Psi(\bs v)+\bs f\cdot(\bs w-\bs v)$ for all $\bs w$.
Taking $\bs w=\lambda\bs v$ and using $\Psi(\lambda\bs v)=\lambda\Psi(\bs v)$,
\[
  \lambda\Psi(\bs v)\ge\Psi(\bs v)+(\lambda-1)\bs f\cdot\bs v
  \ \Longrightarrow\ (\lambda-1)\bigl(\Psi(\bs v)-\bs f\cdot\bs v\bigr)\ge0.
\]
Choosing $\lambda>1$ gives $\Psi(\bs v)\ge\bs f\cdot\bs v$ and
$\lambda\in(0,1)$ gives the reverse inequality, whence
$\bs f\cdot\bs v=\Psi(\bs v)$. Scale invariance of the subdifferential
follows from the identity just proved together with the defining
inequality: for $\lambda>0$, $\bs f\in\partial\Psi(\bs v)$ iff
$\bs f\in\E$ and $\bs f\cdot\bs v=\Psi(\bs v)$, and both conditions are
invariant under $\bs v\mapsto\lambda\bs v$ by $1$-homogeneity. The
stated characterisation is exactly this equivalence, since
$\E=\partial\Psi(\bs 0)=\{\bs f:\bs f\cdot\bs w\le\Psi(\bs w)\ \forall\bs w\}$.
\end{proof}

\begin{proof}[Proof of \cref{prop:biot}]
Fix $t$ where $\dot{\bs q}(t)$ exists. Global stability~(S) implies
local stability, i.e.\ $\bs r=\bs q(t)$ minimises
$\bs r\mapsto E(\bs r,S(t))+\Dr(\bs q(t),\bs r)$; its first-order
optimality condition is $\bs 0\in\partial_{\bs q}E(\bs q,S)+\partial\Psi(\bs 0)$,
that is $\bs f(\bs q,S)\in\E$. Differentiating the energy
balance~(E), which by \cref{prop:power}\ref{E:AC} is absolutely
continuous, gives $\partial_{\bs q}E\cdot\dot{\bs q}+\partial_S E\,\dot S
+\frac{\dd}{\dd t}\Diss_\Psi=0$; since
$\frac{\dd}{\dd t}\Diss_\Psi=\Psi(\dot{\bs q})$ a.e., this reads
$\bs f\cdot\dot{\bs q}=\Psi(\dot{\bs q})$. Combined with $\bs f\in\E$,
\cref{lem:saturation} yields $\bs f\in\partial\Psi(\dot{\bs q})$,
i.e.~\eqref{eq:biot}. Conversely, if $E(\cdot,S)$ is convex, local
stability $\bs f\in\E$ is equivalent to global stability~(S), and
integrating $\bs f\cdot\dot{\bs q}=\Psi(\dot{\bs q})$ recovers~(E).
\end{proof}

% =====================================================================
\section{Proofs of the mathematical results}
\label{app:math}

\begin{proof}[Proof of \cref{prop:quasi_distance}]
Non-negativity is immediate from $\Psi\ge0$, and the constant path
gives $\Dr(\bs q,\bs q)=0$. For the triangle inequality, fix
$\varepsilon>0$ and choose admissible paths $\gamma_{12},\gamma_{23}$
with cost within $\varepsilon/2$ of $\Dr(\bs q_1,\bs q_2)$ and
$\Dr(\bs q_2,\bs q_3)$. Their concatenation, reparametrised to
$[0,1]$, is admissible from $\bs q_1$ to $\bs q_3$; by
$1$-homogeneity the concatenation cost equals the sum of the two
costs, so $\Dr(\bs q_1,\bs q_3)\le\Dr(\bs q_1,\bs q_2)+\Dr(\bs q_2,\bs q_3)+\varepsilon$,
and $\varepsilon\downarrow0$ concludes. If $\Psi$ is symmetric, the
reversed path shows $\Dr(\bs q_2,\bs q_1)=\Dr(\bs q_1,\bs q_2)$; if
moreover $\Psi(\bs v)=0\Rightarrow\bs v=\bs 0$, then $\Dr(\bs q_1,\bs q_2)=0$
forces $\bs q_1=\bs q_2$.
\end{proof}

\begin{proof}[Proof of \cref{prop:Dr_linear}]
For any admissible $\gamma$ from $\bs q_1$ to $\bs q_2$, Jensen's
inequality for the convex, $1$-homogeneous $\Psi$ gives
$\int_0^1\Psi(\dot\gamma)\dd s\ge\Psi\bigl(\int_0^1\dot\gamma\dd s\bigr)
=\Psi(\bs q_2-\bs q_1)$, hence $\Dr\ge\Psi(\bs q_2-\bs q_1)$. Since $\Q$
is convex, the straight segment $\bar\gamma(s)=\bs q_1+s(\bs q_2-\bs q_1)$
is admissible with $\dot{\bar\gamma}\equiv\bs q_2-\bs q_1$ and cost
$\Psi(\bs q_2-\bs q_1)$, giving the reverse inequality. Continuity of
$\Dr$ follows from that of $\Psi$.
\end{proof}

\begin{proof}[Proof of \cref{prop:power}]
\emph{\ref{E:compact}.} If $E(\bs q,S)\le c$ then, by~\ref{A:W}
and~\ref{A:ell} and Cauchy--Schwarz,
$c_1\|\bs q\|^p-c_2\le W(\bs q)\le c+\|\bs q\|C_\ell$; since $p>1$ the
left side dominates and $\|\bs q\|$ is bounded on the sublevel set,
which is closed by continuity of $E(\cdot,S)$, hence compact in
$\R^n$. \emph{\ref{E:AC}.} By~\ref{A:ell} $\bs\ell$ is Lipschitz, hence
differentiable a.e.\ with $\|\bs\ell'\|\le C_{\ell'}$; since
$S\in W^{1,1}$, the chain rule for Sobolev functions gives
$\frac{\dd}{\dd t}\bs\ell(S(t))=\bs\ell'(S(t))\dot S(t)$ a.e., and, $W$
being $t$-independent, $\frac{\dd}{\dd t}E(\bs q,S(t))=-\bs q\cdot\bs\ell'(S)\dot S$.
\emph{\ref{E:gronwall}.} From~\ref{A:W}--\ref{A:ell},
$c_1\|\bs q\|^p\le E(\bs q,S)+C_\ell\|\bs q\|+c_2$; Young's inequality
$C_\ell\|\bs q\|\le\frac{c_1}{2}\|\bs q\|^p+C$ absorbs the linear term,
giving $\|\bs q\|^p\le C'(E(\bs q,S)+c_4)$, and since
$|\partial_S E|=|\bs q\cdot\bs\ell'|\le C_{\ell'}\|\bs q\|$, the bound
$|\partial_S E|\le c_3(E(\bs q,S)+c_4)$ follows.
\end{proof}

\begin{proof}[Proof of \cref{lem:compatibility}]
\emph{(C1).} For a.e.\ $t$, $\partial_S E(\bs q,S(t))=-\bs q\cdot\bs\ell'(S(t))\dot S(t)$
is linear in $\bs q$ and, since $\bs\ell'$ is bounded and
$S\in W^{1,1}\subset C^0$, one has $\bs\ell'(S(t_m))\dot S(t_m)\to\bs\ell'(S(t))\dot S(t)$;
with $\bs q_m\to\bs q$ this gives $\partial_S E(\bs q_m,S(t_m))\to\partial_S E(\bs q,S(t))$.
\emph{(C2).} From $\bs q_m\in\mathcal S(t_m)$,
$E(\bs q_m,S(t_m))\le E(\tilde{\bs q},S(t_m))+\Dr(\bs q_m,\tilde{\bs q})$
for all $\tilde{\bs q}$. Continuity of $W$ and $\bs\ell$
(hence of $E(\cdot,S)$ and, by \cref{prop:Dr_linear} or
\cref{prop:quasi_distance}, of $\Dr$) allows passage to the limit
$m\to\infty$ at fixed $\tilde{\bs q}$, yielding
$E(\bs q,S(t))\le E(\tilde{\bs q},S(t))+\Dr(\bs q,\tilde{\bs q})$, i.e.\
$\bs q\in\mathcal S(t)$.
\end{proof}

\begin{proof}[Proof of \cref{thm:existence}]
We verify the hypotheses of \cite[Thm.~2.1.6]{MielkeRoubicek} for
$(\Q,E,\Dr)$. Conditions (D1)--(D2) hold by
\cref{prop:quasi_distance}, and $\Dr$ is a metric since $\Psi$ is
symmetric. Compactness of energy sublevels and the power control are
\cref{prop:power}\ref{E:compact} and~\ref{E:gronwall}, giving the
abstract conditions on $E$ with $\alpha_E(t)=c_3|\dot S(t)|\in L^1$.
The compatibility conditions (C1)--(C2) are \cref{lem:compatibility}.
Finally $\Q\subset\R^n$ is separable and metrizable. With
$\bs q_0\in\mathcal S(0)$, \cite[Thm.~2.1.6]{MielkeRoubicek} yields an
energetic solution of bounded variation.
\end{proof}

\begin{proof}[Proof of \cref{thm:uniqueness}]
Under uniform convexity of $W$ and linearity of $\bs q\mapsto-\bs q\cdot\bs\ell(S)$,
$E(\cdot,S)$ is strictly convex with $D^2_{\bs q}E\succeq\alpha\Id$, so
the stability condition $\bs q_0\in\mathcal S(0)$ is equivalent to the
first-order inclusion $-\partial_{\bs q}E(\bs q_0,S(0))\in\partial\Psi(\bs 0)$.
Together with the joint $C^3$ regularity of $E$ (from
$\bs\ell,S\in C^3$), these are the hypotheses of
\cite[Thm.~3.4.7]{MielkeRoubicek}, whose Section~3.4.4 gives
uniqueness.
\end{proof}

\begin{proof}[Proof of \cref{thm:BV_existence}]
\emph{Step 1.} For $\varepsilon>0$, \eqref{eq:regularised} is the
gradient flow $\bs 0\in\partial\Phi_\varepsilon(\dot{\bs q}_\varepsilon)+\partial_{\bs q}E$
with $\Phi_\varepsilon(\bs v)=\Psi(\bs v)+\frac{\varepsilon}{2}\|\bs v\|^2$
strictly convex; by maximal-monotone operator theory
\cite{Brezis1973} it has a unique solution
$\bs q_\varepsilon\in W^{1,2}(0,T;\Q)$ satisfying a strict energy
balance. \emph{Step 2.} The Gronwall control
\cref{prop:power}\ref{E:gronwall}, uniform in $\varepsilon$, bounds
$E(\bs q_\varepsilon(t),S(t))\le M$; using~\ref{A:coerc},
$\Psi(\bs v)\ge c_\Psi\|\bs v\|$, the total variation is controlled by
the dissipation,
$\Var(\bs q_\varepsilon;[0,T])\le c_\Psi^{-1}\Diss_\Psi(\bs q_\varepsilon;[0,T])
\le c_\Psi^{-1}[E(\bs q_0,S(0))-E(\bs q_\varepsilon(T),S(T))+\mathcal W(T)]\le C$,
uniformly in $\varepsilon$. \emph{Step 3.} By Helly's selection
theorem \cite[Thm.~3.3]{AmbrosioFuscoPallara} there is a subsequence
$\bs q_{\varepsilon_j}\to\bs q$ pointwise, with $\bs q$ of bounded
variation. Lower semicontinuity of $E(\cdot,S)$ and continuity of
$\Dr$ (\cref{prop:Dr_linear}) preserve (BV-S); the energy balance
passes to the limit by \cref{prop:power}\ref{E:AC}, the viscous
correction $\mathcal V$ being identified through the rescaled jump
analysis as in \cite[Thm.~3.3.2]{MielkeRoubicek}, giving (BV-E).
\end{proof}

\begin{proof}[Proof of \cref{thm:BV_uniqueness}]
\emph{Uniqueness of the post-jump state.} In $n=1$, any monotone path
between $q^-$ and $q^+$ realises the dissipation distance, so the jump
condition of \cref{rem:jump} selects the path of steepest energy
descent through the barrier; since $E(\cdot,S(t^\ast))$ has finitely
many local minima, the first local minimum reached beyond the
losing-stability point is uniquely determined, fixing $q^+(t^\ast)$.
\emph{Uniqueness of the trajectory.} Between consecutive jumps the
local stability (BV-S) and the flow rule $\bs f\in\partial\Psi(\dot q)$
determine $q$ uniquely as the solution of the scalar yield condition
$g(q)=\ell(S)\mp\rho$ on the current branch; concatenating the uniquely
determined smooth arcs and jumps gives a unique BV solution.
\end{proof}

% =====================================================================
\section{Proofs for the variational numerical method}
\label{app:numerical}

\begin{proof}[Proof of \cref{prop:return_map}]
Since $\Dr(\bs q_k,\cdot)$ is convex (\cref{prop:quasi_distance}) and
$E(\cdot,S_{k+1})$ is differentiable, the optimality condition
of~\eqref{eq:Ikstar} is
$\bs 0\in\partial_{\bs q}E(\bs q_{k+1},S_{k+1})+\partial_{\bs q}\Dr(\bs q_k,\bs q_{k+1})$.
By \cref{prop:Dr_linear}, $\partial_{\bs q}\Dr(\bs q_k,\bs q)=\partial\Psi(\bs q-\bs q_k)$,
giving~\eqref{eq:discrete_inclusion}. If $\bs q_{k+1}=\bs q_k$ the
inclusion reduces to $-\partial_{\bs q}E(\bs q_k,S_{k+1})\in\partial\Psi(\bs 0)=\E$;
otherwise $\bs q_{k+1}-\bs q_k\ne\bs 0$ and, by \cref{lem:saturation},
$-\partial_{\bs q}E(\bs q_{k+1},S_{k+1})\in\partial\Psi(\bs q_{k+1}-\bs q_k)\setminus\E$.
In the scalar case, $\partial\Psi(\dot q)=\rho\sgn(\dot q)$ equals
$\{\rho\}$, $[-\rho,\rho]$, $\{-\rho\}$ for $\dot q>0,=0,<0$, and the
inclusion reads $f=\ell(S_{k+1})-g(q_{k+1})\in\rho\sgn(q_{k+1}-q_k)$,
i.e.~\eqref{eq:scalar_return_map}.
\end{proof}

\begin{proof}[Proof of \cref{prop:discrete_energy}]
Minimality of $\bs q_{k+1}$ in~\eqref{eq:Ikstar} gives, for all
$\tilde{\bs q}$,
$E(\bs q_{k+1},S_{k+1})+\Dr(\bs q_k,\bs q_{k+1})\le E(\tilde{\bs q},S_{k+1})+\Dr(\bs q_k,\tilde{\bs q})$;
dropping the non-negative term $\Dr(\bs q_k,\bs q_{k+1})$ on the left
gives~\eqref{eq:disc_stab}. Taking $\tilde{\bs q}=\bs q_k$ and using
$\Dr(\bs q_k,\bs q_k)=0$ gives
$E(\bs q_{k+1},S_{k+1})+\Dr(\bs q_k,\bs q_{k+1})\le E(\bs q_k,S_{k+1})$,
which is~\eqref{eq:disc_energy}.
\end{proof}

\begin{proof}[Proof of \cref{thm:convergence}]
\emph{Uniform bounds.} Summing~\eqref{eq:disc_energy} over
$k=0,\dots,K-1$ and estimating the discrete work by
$|E(\bs q_k,S_{k+1})-E(\bs q_k,S_k)|\le C_{\ell'}\|\bs q_k\|\int_{t_k}^{t_{k+1}}|\dot S|$
(via~\ref{A:ell} and the fundamental theorem of calculus) gives, with
the Gronwall control \cref{prop:power}\ref{E:gronwall} and the discrete
Gronwall lemma, a uniform bound on $E(\bs q_K,S_K)$ and on
$\sum_k\Dr(\bs q_k,\bs q_{k+1})$, hence on
$\Var(\bs q^h;[0,T])$. \emph{Compactness.} By Helly's theorem
\cite[Thm.~3.3]{AmbrosioFuscoPallara} a subsequence $\bs q^{h_j}$
converges pointwise to some $\bs q$ of bounded variation.
\emph{Limit.} The discrete stability~\eqref{eq:disc_stab} passes to the
limit by lower semicontinuity of $E$ and continuity of $\Dr$, giving
(S); the summed energy inequality passes to the limit, using
\cref{prop:power}\ref{E:AC} for the work term and lower semicontinuity
of $\Diss_\Psi$, and the reverse inequality from (S) as in
\cite[Sec.~2.1]{MielkeRoubicek}, giving (E). Thus $\bs q$ is an
energetic solution. Under uniqueness every subsequence has the same
limit, so the whole family converges.
\end{proof}

\begin{proof}[Proof of \cref{prop:consistency}]
By the fundamental theorem of calculus,
$E(\bs q(t_k),S_{k+1})-E(\bs q(t_k),S_k)=\int_{t_k}^{t_{k+1}}\partial_S E(\bs q(t_k),S(s))\dot S(s)\dd s$,
so
\[
  \tau_k=\int_{t_k}^{t_{k+1}}\bigl[\partial_S E(\bs q(s),S(s))-\partial_S E(\bs q(t_k),S(s))\bigr]\dot S(s)\dd s
  =\int_{t_k}^{t_{k+1}}(\bs q(s)-\bs q(t_k))\cdot\bs\ell'(S(s))\,\dot S(s)\dd s,
\]
using $\partial_S E=-\bs q\cdot\bs\ell'(S)$. By Cauchy--Schwarz,
\ref{A:ell} and $S\in W^{1,\infty}$,
$|\tau_k|\le C_{\ell'}\|\dot S\|_{L^\infty}\int_{t_k}^{t_{k+1}}\|\bs q(s)-\bs q(t_k)\|\dd s$.
\emph{(i)} If $\bs q\in\mathrm{AC}$, then
$\|\bs q(s)-\bs q(t_k)\|\le\int_{t_k}^{t_{k+1}}\|\dot{\bs q}\|\dd\tau$;
integrating over $s\in[t_k,t_{k+1}]$ yields~\eqref{eq:cons_weak}.
\emph{(ii)} If $\bs q\in W^{1,\infty}$, then
$\|\bs q(s)-\bs q(t_k)\|\le\|\dot{\bs q}\|_{L^\infty}(s-t_k)$ and
$\int_{t_k}^{t_{k+1}}(s-t_k)\dd s=\tfrac12 h_k^2$,
giving~\eqref{eq:cons_strong}.
\end{proof}

\begin{proof}[Proof of \cref{thm:global_consistency}]
Sum the local bounds. In the weak case,
$\mathcal E_\tau\le C_{\ell'}\|\dot S\|_{L^\infty}\sum_k h_k\int_{t_k}^{t_{k+1}}\|\dot{\bs q}\|\dd\tau
\le C_{\ell'}\|\dot S\|_{L^\infty}h\int_0^T\|\dot{\bs q}\|\dd\tau$.
In the strong case,
$\mathcal E_\tau\le\tfrac12 C_{\ell'}\|\dot{\bs q}\|_{L^\infty}\|\dot S\|_{L^\infty}\sum_k h_k^2
\le\tfrac12 C_{\ell'}\|\dot{\bs q}\|_{L^\infty}\|\dot S\|_{L^\infty}\,T\,h$,
using $\sum_k h_k^2\le h\sum_k h_k=Th$. Both give $\mathcal E_\tau=O(h)$.
\end{proof}


\bibliographystyle{plain}
\bibliography{references}

\end{document}
